\scanpage{263}\mspage{179}
% archive sheet number: 262
conditions impliquent $\deleted{M_\pi}\mathcal F^\pi=M_\pi$, i.e. $M_\pi\subset\mathcal F^\pi$,
dans le cas où $\pi>\pi_0$, se voit comme (4.7.9.1),
en utilisant le lemme 4.7.9.1.

\uline{Définition} 4.7.13 : $1^{\mathrm{o}}$) Un cardinal $\pi\geq\pi_0$ est
dit \inserted{\uline{cardinal}} \uline{utile} pour la fcga $(\mathcal F^\pi)_{\pi\geq\pi_0}$, \deleted{\unsure{si elle|s'il}}
s'il satisfait les conditions équivalentes
a) et b) de 4.7.12, et si de plus, dans
le cas où $\pi=\pi_0$, on a \deleted{(4.7.12.1)} $\mathcal F^\pi=M_\pi$. (\unsure{Donc|Dans}
si $\pi$ est un cardinal utile, on a \inserted{\unsure{tjrs}} \deleted{\unsure{4.7}} $\mathcal F^\pi=M_\pi$.)
\deleted{Donc si $\pi$ est utile, on a}
\deleted{On dit que $\pi$ est un card}
Dans le cas contraire, on dit que $\pi$ est
un \uline{cardinal superflu} (pour la fcga
donnée).

$2^{\mathrm{o}}$) Un cardinal $\pi$ est dit \uline{régulier}
pour la fcg donnée, \deleted{s'il satisfait} \inserted{s'il est utile}
\deleted{\unsure{à}} \inserted{et si de plus} la majoration de cardinaux énoncée dans
(4.7.7.1) est valable.

4.7.\unsure{14}.\deleted{\unsure{1}} Ainsi, tous les cardinaux $\pi\geq\pi_0$ tels que
$\pi^{\pi_0}=\pi$ sont utiles, et même réguliers pour
la filtration cardinale donnée. Notons
que si tous les cardinaux $\pi\geq\pi_0$ sont
utiles, alors on a $\mathcal F^\pi=M_\pi$ pour tout
$\pi\geq\pi_0$, donc la fcga est canonique. Notons
aussi que dans tous les cas, si $(\mathcal F^\pi)_{\pi\geq\pi_0}$ est
une fcga, alors $(M_\pi)_{\pi\geq\pi_0}$ en est une aussi.
